The rapid adoption of Internet of Things (IoT) technologies has
transformed modern computing environments by enabling interconnected
systems across industrial automation, healthcare, smart cities, and
critical infrastructure. However, the expansion of IoT deployments
towards distributed Edge-IoT architectures has significantly increased
the cybersecurity attack surface. Unlike conventional enterprise
networks, Edge-IoT environments consist of heterogeneous devices with
limited computational capabilities, dynamic communication patterns,
and decentralized deployment characteristics.

Traditional intrusion detection systems (IDS) primarily focus on
identifying whether network activity is malicious or benign. Although
machine learning based IDS approaches have demonstrated strong
classification performance, they generally provide limited contextual
understanding regarding attack severity, asset importance, and
operational impact. In practical Security Operations Center (SOC)
environments, analysts require not only detection capability but also
risk prioritization to determine which alerts require immediate
response.

Recent research has explored artificial intelligence techniques for
enhancing IoT security. Dataset-driven approaches such as TON-IoT have
enabled standardized evaluation of intrusion detection models across
heterogeneous IoT and industrial environments. Machine learning models
including Random Forest, Support Vector Machines, and deep learning
architectures have achieved high detection performance; however, these
approaches commonly operate as isolated classifiers without integrated
risk reasoning or response prioritization.

Agentic AI introduces a different perspective by enabling systems to
coordinate multiple decision-making components for autonomous security
operations. In cybersecurity applications, agent-based architectures
can combine threat perception, contextual analysis, risk estimation,
and response planning. However, existing studies frequently focus on
either intrusion detection or automated response, leaving a gap in
integrating accurate detection with adaptive risk assessment for
Edge-IoT environments.

This paper presents RiskAdaptive, an agentic AI framework that extends
traditional intrusion detection by incorporating a risk-aware decision
layer. The proposed system combines machine learning based threat
identification with contextual risk scoring to classify detected
events according to operational severity. Unlike conventional IDS
approaches that terminate after attack classification, RiskAdaptive
provides additional security intelligence for prioritizing incidents
and supporting automated security decision workflows.

The major contributions of this work are summarized as follows:

\begin{itemize}
\item A risk-aware agentic AI framework is proposed for Edge-IoT
intrusion detection and incident prioritization.

\item A dynamic risk assessment layer is integrated with intrusion
classification to transform binary detection outputs into
severity-aware security decisions.

\item Extensive evaluation is performed using the TON-IoT network
telemetry dataset with comparison against conventional machine
learning baselines including Logistic Regression and Random Forest.

\item Attack-specific analysis is presented to evaluate detection
capability and risk distribution across different threat categories.
\end{itemize}

The remainder of this paper is organized as follows. Section II
reviews related research in IoT intrusion detection, agentic AI
security frameworks, and risk-aware cybersecurity approaches. Section
III presents the proposed RiskAdaptive methodology. Section IV
describes the experimental setup and evaluation procedure. Section V
discusses experimental results and findings. Section VI concludes the
paper and highlights future research directions.
